\chapter{引言}

本书介绍线性泛函分析，将经典线性代数的技术与结论推广至无穷维空间.随着函数空间理论的发展，泛函分析成为研究线性常微分方程与偏微分方程的有力工具；它为解的存在性与唯一性、解对初始条件或边界数据的连续依赖性、近似解的收敛性，以及其他各类性质，提供了根本性洞见.

以下评注突出若干线性代数中的关键结果及其在无穷维情形下的对应形式.

\section{线性方程}

设 \(A\) 为一个 \(n \times  n\) 矩阵.给定向量 \(\mathbf{b} \in  {\mathbb{R}}^{n}\) ，线性代数中的一个基本问题是求向量 \(\mathbf{x} \in  {\mathbb{R}}^{n}\) ，使得

\begin{equation}\tag{1.1}
A\mathbf{x} = \mathbf{b}.
\end{equation}


在线性偏微分方程理论中，相应问题如下:考虑一个有界开集 \(\Omega  \subset  {\mathbb{R}}^{n}\) ，以及形如

\begin{equation}\tag{1.2}
{Lu} \doteq   - \mathop{\sum }\limits_{{i,j = 1}}^{n}{\left( {a}^{ij}\left( x\right) {u}_{{x}_{i}}\right) }_{{x}_{j}} + \mathop{\sum }\limits_{{i = 1}}^{n}{b}^{i}\left( x\right) {u}_{{x}_{i}} + c\left( x\right) u.
\end{equation}

的线性偏微分算子.给定函数 \(f : \Omega  \mapsto  \mathbb{R}\) ，求函数 \(u\) ，使其在 \(\Omega\) 的边界上为零，且满足

\begin{equation}\tag{1.3}
{Lu} = f.
\end{equation}

问题(1.1)与(1.3)之间存在本质差异.矩阵 \(A\) 在有限维空间 \({\mathbb{R}}^{n}\) 上诱导一个连续线性变换；而微分算子 \(L\) 则可视为无限维空间 \({\mathbf{L}}^{2}\left( \Omega \right)\) 上的无界(因而不连续)线性算子.特别地， \(L\) 的定义域并非整个空间 \({\mathbf{L}}^{2}\left( \Omega \right)\) ，而仅为其中某个适当的子空间.

尽管存在这些差异，但由于问题(1.1)与(1.3)均为线性问题，故线性代数中诸多方法亦可适用于(1.3).

\section*{(I):正定性}

假设矩阵 \(A\) 严格正定，即存在常数 \(\beta  > 0\) 使得

\[
\langle {Ax},x\rangle  \geq  \beta {\left| x\right| }^{2},\forall x \in  {\mathbb{R}}^{n}.
\]

则 \(A\) 可逆，且对任意 \(\mathbf{b} \in  {\mathbb{R}}^{n}\) ，方程(1.1)均有唯一解.

该结论对椭圆型偏微分方程具有直接对应形式.即假设算子 \(L\) 严格正定，其含义为(经形式分部积分后)

\begin{equation}\tag{1.4}
\begin{aligned}
\langle {Lu},u{\rangle }_{{\mathbf{L}}^{2}} &= {\int }_{\Omega }\left( {\mathop{\sum }\limits_{{i,j = 1}}^{n}{a}^{ij}\left( x\right) {u}_{{x}_{i}}{u}_{{x}_{j}} + \mathop{\sum }\limits_{{i = 1}}^{n}{b}^{i}\left( x\right) {u}_{{x}_{i}}u + \mathrm{c}\left( x\right) {u}^{2}}\right) {dx} \\
&\geq  \beta \| u{\| }_{{H}_{0}^{1}\left( \Omega \right) }^{2}
\end{aligned}
\end{equation}

对某个常数 \(\beta  > 0\) 及所有 \(u \in  {H}_{0}^{1}\left( \Omega \right)\) 成立.此处 \({H}_{0}^{1}\left( \Omega \right)\) 是定义在 \(\Omega\) 边界上取零值、且满足

\[
\| u{\| }_{{H}_{0}^{1}\left( \Omega \right) } \doteq  {\left( {\int }_{\Omega }{\left| u\right| }^{2}dx + {\int }_{\Omega }{\left| \nabla u\right| }^{2}dx\right) }^{1/2} < \infty ;
\]
的函数空间；精确的定义见第8章.若(1.4)成立，则可证问题(1.3)对任意给定 \(f \in  {\mathbf{L}}^{2}\left( \Omega \right)\) 均存在唯一解 \(u \in  {H}_{0}^{1}\left( \Omega \right)\) .

为使不等式(1.4)成立，一个关键假设是算子 \(L\) 应为椭圆型.即在每一点 \(x \in  \Omega\) 处， \(n \times  n\) 阶矩阵 \(\left( {{a}^{ij}\left( x\right) }\right)\) 均须严格正定.

\section*{(II):Fredholm择一性}

线性代数中一个著名判据指出:方程(1.1)对任意给定 \(\mathbf{b} \in  {\mathbb{R}}^{n}\) 均有唯一解，当且仅当齐次方程

\[
A\mathbf{x} = 0
\]

仅有零解 \(\mathbf{x} = 0\) .显然，这成立当且仅当矩阵 \(A\) 可逆.

一般而言，无限维空间 \(X\) 上的连续线性算子并不具备此性质.事实上，可构造出单射但非满射(或反之)的有界线性算子 \(\Lambda  : X \mapsto  X\) .

然而，有限维理论可推广至一类重要算子，即形如 \(\Lambda  = I - K\) 的算子，其中 \(I\) 为恒等算子， \(K\) 为紧算子.若 \(\Lambda\) 属于此类，则仍可证明等价性

\[
\Lambda \text{ is one-to-one } \Leftrightarrow  \Lambda \text{ is onto. }
\]

由此理论的一个应用可知:对线性椭圆型算子，方程(1.3)对任意 \(f \in  {\mathbf{L}}^{2}\left( \Omega \right)\) 均有唯一解 \(u \in  {H}_{0}^{1}\left( \Omega \right)\) ，当且仅当齐次方程

\[
{Lu} = 0
\]

仅有零解.

\section*{(III):对角化}

若能在 \({\mathbb{R}}^{n}\) 中找到由 \(A\) 的特征向量构成的一组基 \(\left\{  {{\mathbf{v}}_{1},\ldots ,{\mathbf{v}}_{n}}\right\}\) ，则系统(1.1)在此基下呈对角形式，从而易于求解.

对于具有重特征值的一般矩阵 \(A\) ，众所周知，这样的特征向量基未必存在.这方面的一个肯定结果如下:若 \(n \times  n\) 矩阵 \(A\) 为对称矩阵，则可在欧氏空间 \({\mathbb{R}}^{n}\) 中找到一组由 \(A\) 的特征向量构成的标准正交基 \(\left\{  {{\mathbf{v}}_{1},\ldots ,{\mathbf{v}}_{n}}\right\}\) .即:

\[
\left( {{\mathbf{v}}_{i},{\mathbf{v}}_{j}}\right)  = \left\{  \begin{array}{ll} 1 & \text{ if }i = j, \\  0 & \text{ if }i \neq  j, \end{array}\right. \;A{\mathbf{v}}_{k} = {\lambda }_{k}{\mathbf{v}}_{k}.
\]

此处 \({\lambda }_{1},\ldots ,{\lambda }_{n} \in  \mathbb{R}\) 为对应的特征值.现可通过计算 \(\mathbf{x}\) 在该标准正交基下的系数 \({c}_{1},\ldots ,{c}_{n}\) ，求得(1.1)的解:

\[
\mathbf{x} = \mathop{\sum }\limits_{{k = 1}}^{n}{c}_{k}{\mathbf{v}}_{k},\;A\mathbf{x} = \mathop{\sum }\limits_{{k = 1}}^{n}{\lambda }_{k}{c}_{k}{\mathbf{v}}_{k} = \mathbf{b} = \mathop{\sum }\limits_{{k = 1}}^{n}\left\langle  {\mathbf{b},{\mathbf{v}}_{k}}\right\rangle  {\mathbf{v}}_{k}.
\]

注意，得益于特征向量基，问题得以解耦:无需求解一个含 \(n\) 个方程、 \(n\) 个未知量的大系统，而只需分别求解 \(n\) 个标量方程，每个对应一个系数 \({c}_{k}\) .若所有特征值 \({\lambda }_{k}\) 均非零，则可得到显式公式:

\begin{equation}\tag{1.5}
\mathbf{x} = \mathop{\sum }\limits_{{k = 1}}^{n}\frac{1}{{\lambda }_{k}}\left\langle  {\mathbf{b},{\mathbf{v}}_{k}}\right\rangle  {\mathbf{v}}_{k}.
\end{equation}

在分析(1.2)中的椭圆型算子 \(L\) 时，只要满足 \({a}^{ij} = {a}^{ji}\) 和 \({b}^{i}\left( x\right)  = 0\) ，即可采用相同方法.事实上，这些条件使该算子具有“对称性”.此时可在空间 \({\mathbf{L}}^{2}\left( \Omega \right)\) 中找到一组可数的标准正交基 \(\left\{  {{\phi }_{1},{\phi }_{2},\ldots }\right\}\) ，其元素为函数 \({\phi }_{k} \in  {H}_{0}^{1}\left( \Omega \right)\) ，使得

\begin{equation}\tag{1.6}
{\left\langle  {\phi }_{i},{\phi }_{j}\right\rangle  }_{{\mathbf{L}}^{2}} = \left\{  \begin{array}{ll} 1 & \text{ if }i = j, \\  0 & \text{ if }i \neq  j, \end{array}\right. \;L{\phi }_{k} = {\lambda }_{k}{\phi }_{k},
\end{equation}

其中 \({\lambda }_{k} \rightarrow   + \infty\) 为一适当选取的实特征值序列.若假设对所有 \(k\) 均有 \({\lambda }_{k} \neq  0\) ，则(1.3)的唯一解可显式写为:

\begin{equation}\tag{1.7}
u = \mathop{\sum }\limits_{{k = 1}}^{\infty }\frac{1}{{\lambda }_{k}}{\left\langle  f,{\phi }_{k}\right\rangle  }_{{\mathbf{L}}^{2}}{\phi }_{k}.
\end{equation}

注意公式(1.5)与(1.7)之间的高度相似性.本质上，只需将 \({\mathbb{R}}^{n}\) 上的欧氏内积替换为 \({\mathbf{L}}^{2}\left( \Omega \right)\) 上的内积即可.

\section{演化方程}

设 \(A\) 为一个 \(n \times  n\) 矩阵.对给定初态 \(\mathbf{b} \in  {\mathbb{R}}^{n}\) ，考虑柯西问题:

\begin{equation}\tag{1.8}
\frac{d}{dt}\mathbf{x}\left( t\right)  = A\mathbf{x}\left( t\right) ,\;\mathbf{x}\left( 0\right)  = \mathbf{b}.
\end{equation}

根据线性常微分方程理论，该问题存在唯一解:

\begin{equation}\tag{1.9}
\mathbf{x}\left( t\right)  = {e}^{tA}\mathbf{b},
\end{equation}

其中

\begin{equation}\tag{1.10}
{e}^{tA} = \mathop{\sum }\limits_{{k = 0}}^{\infty }\frac{{t}^{k}{A}^{k}}{k!}
\end{equation}

注意，(1.10)式右端被定义为一个收敛的 \(n \times  n\) 矩阵级数；此处 \({A}^{0} = I\) 为单位矩阵.矩阵族 \(\left\{  {{e}^{tA};t \in  \mathbb{R}}\right\}\) 具有“群性质”，即:

\[
{e}^{0A} = I,\;{e}^{tA}{e}^{sA} = {e}^{\left( {t + s}\right) A}\;\text{ for all }t,s \in  \mathbb{R}.
\]

若 \(A\) 为对称矩阵，则它存在一组标准正交特征向量基 \(\left\{  {{\mathbf{v}}_{1},\ldots ,{\mathbf{v}}_{n}}\right\}\) ，对应特征值为 \({\lambda }_{1},\ldots ,{\lambda }_{n}\) .此时，解(1.9)可更显式地写为:

\[
{e}^{tA}\mathbf{b} = \mathop{\sum }\limits_{{k = 1}}^{n}{e}^{t{\lambda }_{k}}\left\langle  {\mathbf{b},{\mathbf{v}}_{k}}\right\rangle  {\mathbf{v}}_{k}.
\]

线性半群理论将上述结果推广至无穷维空间中的无界线性算子.特别地，它适用于如下形式的抛物型演化方程:

\begin{equation}\tag{1.11}
\frac{d}{dt}u\left( t\right)  =  - {Lu}\left( t\right) ,\;u\left( 0\right)  = g \in  {\mathbf{L}}^{2}\left( \Omega \right) ,\;u = 0\text{ on }\partial \Omega ,
\end{equation}

其中 \(L\) 是 (1.2) 式中的偏微分算子， \(\partial \Omega\) 表示 \(\Omega\) 的边界.当 \({a}^{ij} = {a}^{ji}\) 且 \({b}^{i}\left( x\right)  = 0\) 时，椭圆算子 \(L\) 是对称的，其解可沿 (1.6) 式中考虑的空间 \({\mathbf{L}}^{2}\left( \Omega \right)\) 的标准正交基 \(\left\{  {{\phi }_{1},{\phi }_{2},\ldots }\right\}\) 分解.由此得到如下表示

\begin{equation}\tag{1.12}
u\left( t\right)  = {S}_{t}g \doteq  \mathop{\sum }\limits_{{k = 1}}^{\infty }{e}^{-t{\lambda }_{k}}{\left\langle  g,{\phi }_{k}\right\rangle  }_{{\mathbf{L}}^{2}}{\phi }_{k},\;t \geq  0.
\end{equation}

注意，算子 \(L\) 是无界的(其特征值满足 \({\lambda }_{k} \rightarrow   + \infty\) ，当 \(k \rightarrow  \infty\) 时).然而，(1.12) 式中的算子 \({S}_{t}\) 对每个 \(t \geq  0\) 均有界(但对 \(t < 0\) 不成立).线性算子族 \(\left\{  {{S}_{t};t \geq  0}\right\}\) 称为线性半群，因其具有半群性质

\[
{S}_{0} = I,\;{S}_{t} \circ  {S}_{s} = {S}_{t + s}\;\text{ for all }s,t \geq  0.
\]

直观上，可将 \({S}_{t}\) 视为指数型算子: \({S}_{t} \doteq  {e}^{-{Lt}}\) .但由于 \(L\) 是无界的，须注意类似 (1.10) 式的指数公式不再成立.当显式公式 (1.12) 不可用时，必须借助其他近似方法构造算子 \({S}_{t}\) .在有限维情形下，矩阵 \(A\) 的指数可通过下式复原

\begin{equation}\tag{1.13}
{e}^{tA} = \mathop{\lim }\limits_{{n \rightarrow  \infty }}{\left( I - \frac{t}{n}A\right) }^{-n}
\end{equation}

也可通过下式复原

\begin{equation}\tag{1.14}
{e}^{tA} = \mathop{\lim }\limits_{{\lambda  \rightarrow  \infty }}{e}^{t{A}_{\lambda }},\;{A}_{\lambda } \doteq  A{\left( I - {\lambda }^{-1}A\right) }^{-1}.
\end{equation}

值得注意的是，两个公式 (1.13)-(1.14) 对一大类无穷维空间上的无界算子仍保持有效性.

双曲型初值问题

\begin{equation}\tag{1.15}
{u}_{tt} + {Lu} = 0,\;\begin{cases} u\left( 0\right) &  = g, \\  {u}_{t}\left( 0\right) &  = h, \end{cases}\;u = 0\text{ on }\partial \Omega ,
\end{equation}

亦可用类似方法处理.

(1.15) 式的有限维对应形式是二阶线性方程组

\begin{equation}\tag{1.16}
\frac{{d}^{2}}{d{t}^{2}}\mathbf{x}\left( t\right)  + A\mathbf{x}\left( t\right)  = 0,\;\mathbf{x}\left( 0\right)  = \mathbf{a},\;\frac{d}{dt}\mathbf{x}\left( 0\right)  = \mathbf{b}.
\end{equation}

此处 \(\mathbf{x},\mathbf{a},\mathbf{b} \in  {\mathbb{R}}^{n}\) 和 \(A\) 是一个 \(n \times  n\) 矩阵.用上点号表示时间导数，并令 \(\mathbf{y} \doteq  \dot{\mathbf{x}},\left( {1.16}\right)\) ，则可将 (1.15) 式写成一阶系统:

\begin{equation}\tag{1.17}
\left( \begin{array}{l} \dot{\mathbf{x}} \\  \dot{\mathbf{y}} \end{array}\right)  = \left( \begin{matrix} 0 & I \\   - A & 0 \end{matrix}\right) \left( \begin{array}{l} \mathbf{x} \\  \mathbf{y} \end{array}\right) ,\;\left( \begin{array}{l} \mathbf{x}\left( 0\right) \\  \mathbf{y}\left( 0\right)  \end{array}\right)  = \left( \begin{array}{l} \mathbf{a} \\  \mathbf{b} \end{array}\right) .
\end{equation}

因此，适用于一阶线性常微分方程的结果亦可在此应用.若 \(A\) 对称，则它具有一组标准正交特征向量基 \(\left\{  {{\mathbf{v}}_{1},\ldots ,{\mathbf{v}}_{n}}\right\}\) ，对应特征值为 \({\lambda }_{1},\ldots ,{\lambda }_{n}\) .此时，(1.16) 式的解可表为

\begin{equation}\tag{1.18}
\mathbf{x}\left( t\right)  = \mathop{\sum }\limits_{{k = 1}}^{n}{c}_{k}\left( t\right) {\mathbf{v}}_{k}.
\end{equation}

每个系数 \({c}_{k}\left( \cdot \right)\) 均可独立计算，即求解如下二阶标量常微分方程

\[
\frac{{d}^{2}}{d{t}^{2}}{c}_{k}\left( t\right)  + {\lambda }_{k}{c}_{k}\left( t\right)  = 0,\;{c}_{k}\left( 0\right)  = \left\langle  {\mathbf{a},{v}_{k}}\right\rangle  ,\;\frac{d}{dt}{c}_{k}\left( 0\right)  = \left\langle  {\mathbf{b},{v}_{k}}\right\rangle  .
\]

回到问题 (1.15)，若椭圆算子 \(L\) 对称，则解仍可沿 (1.6) 式中考虑的空间 \({\mathbf{L}}^{2}\left( \Omega \right)\) 的标准正交基 \(\left\{  {{\phi }_{1},{\phi }_{2},\ldots }\right\}\) 分解.由此得到完全类似的表示

\begin{equation}\tag{1.19}
u\left( t\right)  = \mathop{\sum }\limits_{{k = 1}}^{\infty }{c}_{k}\left( t\right) {\phi }_{k}\;t \geq  0,
\end{equation}

其中每个函数 \({c}_{k}\) 由如下方程确定

\[
\frac{{d}^{2}}{d{t}^{2}}{c}_{k}\left( t\right)  + {\lambda }_{k}{c}_{k}\left( t\right)  = 0,\;{c}_{k}\left( 0\right)  = {\left\langle  g,{\phi }_{k}\right\rangle  }_{{\mathbf{L}}^{2}},\;\frac{d}{dt}{c}_{k}\left( 0\right)  = {\left\langle  h,{\phi }_{k}\right\rangle  }_{{\mathbf{L}}^{2}}.
\]

\section{函数空间}

在泛函分析中，一个核心思想是将函数 \(f : {\mathbb{R}}^{n} \mapsto  \mathbb{R}\) 视为抽象向量空间中的点.关于函数的全部信息被浓缩为一个单一数值 \(\| f\|\) ，我们称之为 \(f\) 的范数.通常，范数衡量 \(f\) 及其直至某阶 \(k\) 的偏导数的“大小”.仅依赖这一概念并结合向量空间结构，竟能获得如此众多的结果，实属惊人.这正是泛函分析方法广获成功的原因所在.

将泛函分析应用于积分方程或微分方程时，需建立函数空间理论.在此方向上，自然考虑函数空间 \({\mathcal{C}}^{k}\) ，即所有直至 \(k\) 阶偏导数均有界且连续的函数构成的空间.函数 \(f \in  {\mathcal{C}}^{k}\left( {\mathbb{R}}^{n}\right)\) 的“大小”由如下范数刻画

\[
\| f{\| }_{{\mathcal{C}}^{k}} \doteq  \mathop{\max }\limits_{{{\alpha }_{1} + \cdots  + {\alpha }_{n} \leq  k}}\mathop{\sup }\limits_{{x \in  {\mathbb{R}}^{n}}}\left| {{\partial }_{{x}_{1}}^{{\alpha }_{1}}\cdots {\partial }_{{x}_{n}}^{{\alpha }_{n}}f\left( x\right) }\right| .
\]

然而，空间 \({\mathcal{C}}^{k}\) 并不总适用于偏微分方程的研究.事实上，基于物理或几何考虑，人们往往无法给出解及其导数的最大值估计，而只能给出其 \({\mathbf{L}}^{p}\) 范数的估计，其中 \(p \geq  1\) 为某给定值.这促使人们引入S o b lev空间 \({W}^{k,p}\) ，它包含所有直至 \(k\) 阶导数均属于 \({\mathbf{L}}^{p}\) 的函数.此时，函数 \(f \in  {W}^{k,p}\left( {\mathbb{R}}^{n}\right)\) 的“大小”由如下范数刻画

\[
\| f{\| }_{{W}^{k,p}} \doteq  {\left( \mathop{\sum }\limits_{{{\alpha }_{1} + \cdots  + {\alpha }_{n} \leq  k}}{\int }_{{\mathbb{R}}^{n}}{\left| {\partial }_{{x}_{1}}^{{\alpha }_{1}}\cdots {\partial }_{{x}_{n}}^{{\alpha }_{n}}f\left( x\right) \right| }^{p}dx\right) }^{1/p}
\]

鉴于其在偏微分方程理论中的基础性作用，第8章将全部用于研究Sobolev空间.

\section{紧性}

求解方程时，若无法给出解的显式表达式，一种常用方法依赖于以下三个步骤:

(i) 构造一列近似解 \({\left( {u}_{n}\right) }_{n \geq  1}\) .

(ii) 提取一列收敛子列 \({u}_{{n}_{j}} \rightarrow  \bar{u}\) .

(iii) 证明该极限 \(\bar{u}\) 是原方程的解.

当进行到第(ii)步时，会遇到Euclid空间 \({\mathbb{R}}^{N}\) 与抽象函数空间之间的一个重要差异:即在 \({\mathbb{R}}^{N}\) 中，所有闭有界集都是紧的；换言之，在 \({\mathbb{R}}^{N}\) 中Bolzano-Weierstrass定理成立:

\begin{itemize}\relax
\item\relax 对任意有界序列 \({\left( {u}_{n}\right) }_{n \geq  1}\) ，均可提取出一个收敛子列.
\end{itemize}

如第2章所证(见定理2.22)，这一关键性质在任一有限维赋范空间中均成立，但在任一无限维赋范空间中均不成立.在函数空间中，仅证明近似解序列有界(即对某个常数 \(C\) 及所有 \(n \geq  1\) 满足 \(\|{u}_{n}\| \leq  C\) )并不能保证存在收敛子列.为克服这一根本困难，可采用两种主要途径.

(i) 引入一种更弱的收敛概念，证明每个有界序列(即使在无限维空间中)均存在一个在此弱意义下收敛的子列.

朝此方向的一个关键结果是Banach-Alaoglu定理，该定理将在第2章末尾给出证明；Hilbert空间中的弱收敛则在第5章中讨论.

(ii) 考虑两个不同的范数，记为 \(\| u{\| }_{\text{ weak }} \leq  \| u{\| }_{\text{ strong }}\) ，并满足如下性质:若序列 \({\left( {u}_{n}\right) }_{n \geq  1}\) 在强范数下有界(即 \({\|{u}_{n}\|}_{\text{ strong }} \leq  C\) )，则存在一个子列在弱范数下收敛: \({\|{u}_{{n}_{j}} - \bar{u}\|}_{\text{ weak }} \rightarrow  0\) ，其中极限为某个 \(\bar{u}\) .第3章证明的Ascoli定理及第8章证明的雷利希-孔德拉绍夫紧嵌入定理，分别提供了可实施该方法的不同场景.

偏微分方程分析的很大一部分最终依赖于先验估计的推导.问题本身的性质决定了所能期望的先验界类型，从而也决定了解所在的具体函数空间.这正解释了当前文献中出现多种函数空间的原因.

尽管泛函分析技术具有高度一般性，并能以直观而经济的方式得出基础性结论，但需注意:仅靠泛函分析方法尚不能涵盖偏微分方程理论的全部内容.通常，这些抽象方法所构造的解位于某Sobolev空间中，其函数仅具备使方程有意义所需的最低正则性.对于若干椭圆型与抛物型方程，已知其解实际上具有高得多的正则性；但此类正则性唯有通过更细致的分析才能确立.其他性质——例如椭圆型或抛物型方程的最大值原理、双曲型方程的有限传播速度——同样需要专为偏微分方程设计的额外技巧.对于这些超出本讲义范围的问题，我们参阅专著[E, GT, McO, P, PW, RR, S, T].
